\section*{Appendices: Reusable Field Instruments}

\subsection*{Appendix A: The One-Page Differentiation Thesis}

Complete this template for every serious customer/offer hypothesis. Treat blank fields as unverified assumptions:

\begin{table*}[t]
\centering
\small
\begin{tabularx}{\textwidth}{>{\raggedright\arraybackslash}p{5.5cm} >{\raggedright\arraybackslash}X >{\raggedright\arraybackslash}p{4cm}}
\toprule
\textbf{Field} & \textbf{Working Answer} & \textbf{Evidence Status} \\
\midrule
Target buyer and triggering situation & Working software engineer facing production AI deployment & Observed in 15 interviews \\
Job-to-be-done and cost of inaction & Build reliable, defensible systems; avoid architectural rework & Reported by practitioners \\
Actual alternatives and status quo & Free tutorials, ad-hoc coworker reviews, trial-and-error & Observed behavior \\
Target segment's distinctive priorities & Demands expert critique; accepts rigorous schedule constraints & Inferred from trade-off study \\
Category frame buyers already understand & Advanced applied-AI technical studio & Tested for comprehension \\
Points of parity needed to qualify & Modern tooling, clear curriculum, competent instructor & Table stakes \\
Proposed difference that changes choice & Supervised architecture reviews and oral technical defense & Core differentiator \\
Operating system making it real & Calibrated rubrics, senior reviewers, strict capacity caps & Internal operating rules \\
Evidence visible before purchase & Published evaluation rubrics, sample teardowns, instructor code & Publicly verifiable \\
Outcome observable after purchase & Deployable system with documented latency, drift, and failure tests & Assessed work artifact \\
Deliberate exclusions and trade-offs & No beginner coding; no passive recordings; no placement guarantees & Strategic refusal \\
Acquisition route and buying stakeholders & Technical blogs, GitHub teardowns, engineering community referrals & Low-cost direct channel \\
Price, capacity, and contribution & ₹35,000; 12 learners per cohort; 40\% net contribution & Unit economic model \\
Competitor response and imitation risk & Rivals copy marketing headlines; unable to replicate reviewer depth & Structural friction \\
Accumulating asset & Calibrated benchmark repository and reputation for rigor & Compounding asset \\
Cheapest honest paid test & 3-week diagnostic sprint with paid benchmark review & Initial pilot experiment \\
Falsification condition and review date & 60\% of prospects refuse price; pilot drops below 80\% completion & Explicit stop metric \\
\bottomrule
\end{tabularx}
\caption{The one-page differentiation thesis template.}
\label{tab:ch8-app-thesis}
\end{table*}

\subsection*{Appendix B: The Market-Frame Comparison Card}

Test candidate category frames (such as ``AI technical course,'' ``applied-AI design studio,'' or ``technical capability diagnostic'') by posing 8 questions to qualified target buyers:

\begin{enumerate}
    \item What do you think this offering actually does based on the title alone?
    \item Who do you believe this offering is intended for?
    \item What would you expect this to cost, and what reference product are you comparing it to?
    \item Which 2 real alternatives would you compare this against if you needed this outcome?
    \item Which specific capabilities or features would you assume are included by default?
    \item What evidence would you demand to see before committing your own funds?
    \item What would leave you feeling disappointed or misled after purchasing?
    \item Which words in the description cause confusion, skepticism, or unrealistic expectations?
\end{enumerate}

\subsection*{Appendix C: The Proof Ladder for a New Organization}

Match marketing claims to the proof assets currently accumulated:

\begin{table*}[t]
\centering
\small
\begin{tabularx}{\textwidth}{>{\raggedright\arraybackslash}p{3cm} >{\raggedright\arraybackslash}X >{\raggedright\arraybackslash}X}
\toprule
\textbf{Organizational Stage} & \textbf{Permissible Proof Assets} & \textbf{Mandatory Caveat} \\
\midrule
Day zero & Founder-created demonstrations, technical white papers, public code repositories, transparent rubrics & Proves individual competence, not institutional track record \\
First pilot & Paid transactions, documented delivery hours, baseline and final capability assessments & Small, self-selected sample; limited statistical generalizability \\
Several repeat cohorts & Stable completion rates, rubric reliability data, verified learner testimonials, organic referrals & Cohort selection and founder enthusiasm may still bias outcomes \\
Established programs & Multi-instructor grading calibration, independent advisory board reviews, verified employer feedback & Does not equal accredited university degree status \\
Mature institution & Peer-reviewed research output, external accreditation, longitudinal alumni career tracking & Institutional legitimacy must be continually defended and renewed \\
\bottomrule
\end{tabularx}
\caption{The proof ladder across organizational maturity stages.}
\label{tab:ch8-app-proof-ladder}
\end{table*}

\subsection*{Appendix D: An Activity-System Worksheet}

Map how core positioning promises connect to operational activities and constraints:

\begin{table*}[t]
\centering
\small
\begin{tabularx}{\textwidth}{>{\raggedright\arraybackslash}p{3.5cm} >{\raggedright\arraybackslash}X >{\raggedright\arraybackslash}p{3cm} >{\raggedright\arraybackslash}X}
\toprule
\textbf{Position Promise} & \textbf{Required Operational Activity} & \textbf{Constraint / Cost} & \textbf{Complementary Activity} \\
\midrule
Authentic technical feedback & Calibrated practitioner reviews using published, strict rubrics & Scarce expert hours; review turnaround limits & Selective admissions prerequisites to ensure code readiness \\
Defensible independent work & Mandatory oral technical defense and live failure-mode probing & 45 minutes of evaluator time per learner & Capstone project based on real, non-trivial deployment problems \\
Working-professional fit & Structured, predictable weekly milestone deadlines & Inflexible calendar; no infinite self-pacing & Recorded foundational concept videos and asynchronous chat \\
Evidence-led credibility & Transparent publication of baseline versus final cohort metrics & Privacy governance and assessment administration & Explicit refusal to promise placement or salary outcomes \\
\bottomrule
\end{tabularx}
\caption{Activity-system worksheet demonstrating operational fit and trade-offs.}
\label{tab:ch8-app-activity-system}
\end{table*}

\subsection*{Appendix E: The Perceptual-Map Research Specification}

When conducting perceptual mapping research:
\begin{itemize}
    \item \textbf{Target population:} Define exact screening criteria (such as software engineers with 3 or more years of experience who have deployed at least 1 production service).
    \item \textbf{Familiarity threshold:} Include only alternatives that respondents have actually evaluated, used, or seriously researched.
    \item \textbf{Dimension derivation:} Disclose whether axes represent researcher-selected attributes or statistical factors derived from multi-dimensional scaling.
    \item \textbf{Price anchoring:} Display realistic price points alongside feature descriptions to prevent respondents from rating luxury configurations as universally desirable.
    \item \textbf{Behavioral cross-check:} Pair perceptual ratings with forced-choice trade-off experiments or real deposit commitments.
\end{itemize}

\subsection*{Appendix F: The Positioning Experiment Register}

Track positioning hypotheses systematically:

\begin{table*}[t]
\centering
\small
\begin{tabularx}{\textwidth}{>{\raggedright\arraybackslash}p{2.5cm} >{\raggedright\arraybackslash}X >{\raggedright\arraybackslash}X >{\raggedright\arraybackslash}p{3cm} >{\raggedright\arraybackslash}X}
\toprule
\textbf{Hypothesis} & \textbf{Target Buyer} & \textbf{Offer vs. Alternative} & \textbf{Costly Commitment} & \textbf{Observed Result \& Decision} \\
\midrule
Review intensity outweighs flexibility & Senior software engineers & ₹35,000 reviewed studio vs. free self-study & Paid deposit and 10 hours/week commitment & Recorded in cohort 1; adjust pricing and capacity \\
Diagnostic sprint lowers entry friction & Mid-level ML engineers & ₹10,000 architecture audit vs. full course & Upfront payment and code submission & Test conversion rate to flagship program \\
\bottomrule
\end{tabularx}
\caption{Positioning experiment register for empirical tracking.}
\label{tab:ch8-app-experiment-register}
\end{table*}

\subsection*{Appendix G: Fifteen Questions for Actual Customers}

Administer these chronological questions during discovery and win--loss interviews:

\begin{enumerate}
    \item Walk me through the last time you faced this specific technical or business problem.
    \item What exact event or trigger prompted you to actively search for a solution at that moment?
    \item What was your default approach before you began looking at external alternatives?
    \item Which specific products, services, or internal workarounds did you seriously evaluate?
    \item What initial assumptions or reservations did you hold regarding our offering?
    \item What ultimately made you choose this option over the other alternatives on your shortlist?
    \item What significant compromise, downside, or extra effort did you accept in making that choice?
    \item What feature or service did suppliers pitch that you found entirely irrelevant to your decision?
    \item What proof or evidence was decisive in convincing you that the promise was genuine?
    \item Who else had to approve this expenditure, and what were their primary objections?
    \item What turned out to be unexpectedly difficult or frustrating after you started using the product?
    \item Did the realized outcome match what you believed you were purchasing?
    \item If this solution were discontinued tomorrow, what specific alternative would you adopt?
    \item Would you make the same choice again today, knowing what you now know?
    \item Can you show me a non-confidential artifact or workflow illustrating how you use this solution?
\end{enumerate}

\subsection*{Appendix H: The First 30 Days: A Differentiation Research Sprint}

Execute a structured 4-week differentiation sprint:
\begin{itemize}
    \item \textbf{Week 1: Map the alternatives.} Interview 10 to 12 recent buyers and non-buyers in a single target segment. Document the default habits, costs of inaction, and switching friction.
    \item \textbf{Week 2: Design competing propositions.} Formulate 3 distinct value propositions based on observed customer frustrations. Detail the operational activities and trade-offs required by each.
    \item \textbf{Week 3: Test comprehension and commitment.} Present realistic offers with transparent pricing to qualified prospects. Observe comprehension, measure willingness to pay, and record objections verbatim.
    \item \textbf{Week 4: Deliver and evaluate.} Deliver the smallest viable paid experience to early adopters. Measure baseline versus final capabilities, reviewer hours, net contribution, and customer satisfaction.
\end{itemize}

\subsection*{Appendix I: Positioning Governance for Multiple Products}

Before launching a secondary product under an existing institutional brand, answer 5 governance questions:
\begin{enumerate}
    \item Does the new offering adhere to the master brand's core standards of quality, rigor, and honesty?
    \item Does the target buyer understand how this product differs from existing portfolio offerings?
    \item Does the new product leverage shared operational capabilities, or does it fragment internal resources?
    \item Does this offering introduce reputational, legal, or regulatory risks that could contaminate the core brand?
    \item Does the brand architecture clarify buyer choice, or does it multiply market confusion?
\end{enumerate}

\subsection*{Appendix J: The Founder's Anti-Positioning List}

Establish explicit boundaries by publishing an anti-positioning statement for the organization:
\begin{quote}
\textbf{What Bayes Institute Is Not:}
\begin{itemize}
    \item We are not a vocational placement agency; we do not sell job guarantees or salary promises.
    \item We are not an accredited degree-granting university; we do not award academic degrees.
    \item We are not a passive video library; we do not offer self-paced certificates without human review.
    \item We are not a coding bootcamp for beginners; we do not teach introductory programming syntax.
    \item We are not an on-demand tutoring service; we do not provide unlimited, unstructured founder access.
    \item We are not an AI-shortcut generator; we do not teach students to outsource critical thinking to prompt engines.
\end{itemize}
\end{quote}

Every item on this list protects operational integrity, reduces buyer misalignment, and preserves institutional credibility.
